\documentclass[a4paper,12pt]{article}
\usepackage[utf8]{inputenc}
\usepackage{amsmath,amssymb}
\usepackage{array}
\usepackage{booktabs}
\usepackage{geometry}
\usepackage{parskip}
\usepackage{enumitem}
\usepackage{caption}
\usepackage{hyperref}
\usepackage{tabularx}
\usepackage{multirow}
\usepackage{xcolor}
\usepackage{tcolorbox}

\geometry{margin=2.5cm}

\newtcolorbox{notebox}[1][]{colback=blue!5!white, colframe=blue!50!black, fonttitle=\bfseries, title=#1}
\newtcolorbox{warnbox}[1][]{colback=orange!5!white, colframe=orange!50!black, fonttitle=\bfseries, title=#1}

\title{\textbf{RSE3010: Mine Geotechnical Engineering} \\ Week 7 Study Notes}
\author{Monash University, Semester 1, 2026}
\date{Based on Week 7 Lecture Slides}

\begin{document}

\maketitle

\section{Week 7: Analytical Methods -- Convergence-Confinement Method (CCM)}

\subsection{Purpose}
To introduce analytical solutions for stress and deformation around tunnels and to apply the Convergence-Confinement Method (CCM) for preliminary underground support design. This week integrates concepts from previous weeks:
\begin{itemize}
    \item material behaviour and parameter derivation (Weeks 2--4),
    \item reliability and uncertainty of investigation inputs (Week 5),
    \item failure mechanisms and empirical support design (Week 6),
\end{itemize}
and prepares for numerical modelling (FEM) in Week 8.

\subsection{Learning Objectives}
\begin{itemize}
    \item Apply Kirsch (1898) elastic solutions to estimate stress concentration around a circular excavation.
    \item Understand elastic--plastic zoning and estimate the plastic zone radius using Hoek (1998) (Mohr--Coulomb).
    \item Construct and interpret the Ground Reaction Curve (GRC).
    \item Construct and interpret the Support Characteristic Curve (SCC).
    \item Use a Longitudinal Deformation Profile (LDP) to account for excavation advance and support timing.
    \item Determine the equilibrium pressure $p_{eq}$ and compute the factor of safety, $\text{FoS} = p_{sm}/p_{eq}$.
    \item Recognise the assumptions and limitations of analytical methods and the CCM framework.
\end{itemize}

\subsection{Assumptions and Limitations of Analytical Methods}
All analytical methods presented this week rely on a set of common assumptions that must be understood prior to application:

\begin{itemize}
    \item \textbf{Circular tunnel geometry} -- the solutions are strictly valid only for circular openings.
    \item \textbf{Plane strain conditions} -- the tunnel is considered infinitely long, so no end effects are included.
    \item \textbf{Hydrostatic or simplified stress fields} -- most closed-form solutions assume a hydrostatic far-field stress ($K=1$) or axisymmetric loading.
    \item \textbf{Idealised material behaviour} -- the rock mass is modelled as elastic, elastic-perfectly plastic, or using simplified constitutive laws.
    \item \textbf{Isotropic, homogeneous rock mass} -- anisotropy and large-scale geological structures are neglected.
    \item \textbf{No explicit representation of joints or blocks} -- structurally controlled failures require different analytical or numerical tools.
\end{itemize}

These methods provide first-order estimates and valuable design insights, but they must always be validated against site observations and, where appropriate, supplemented by numerical modelling (Week 8).

\subsection{Elastic Solutions -- Kirsch (1898)}

\subsubsection{Purpose of Kirsch Solutions}
The Kirsch solutions describe the complete stress and displacement fields around a circular opening excavated in a linear-elastic, homogeneous, isotropic medium subjected to biaxial far-field stresses. They form the basis for estimating elastic stress concentrations and serve as the starting point for more advanced elastic--plastic analyses.

\subsubsection{Boundary Stresses at the Tunnel Wall ($r = a$)}
At the tunnel wall, the radial and shear stresses vanish, and the tangential stress is given by:
\[
\sigma_{\theta\theta} = p\left[(1 + K) + 2(1 - K)\cos 2\theta\right],
\]
where $p$ is the vertical far-field stress, $K = \sigma_h / \sigma_v$ is the horizontal-to-vertical stress ratio, and $\theta$ is the angular coordinate measured from the horizontal axis.

\textbf{Critical locations:}
\begin{itemize}
    \item Crown and invert ($\theta = 90^\circ$, $270^\circ$): $\sigma_{\theta\theta} = p(3K - 1)$.
    \item Side walls ($\theta = 0^\circ$, $180^\circ$): $\sigma_{\theta\theta} = p(3 - K)$.
\end{itemize}

\textbf{Sign convention:} Compressive stresses are taken as positive; negative values indicate tension.

\subsubsection{Engineering Applications}
\begin{itemize}
    \item Compare the tangential boundary stress with the rock mass strength (e.g., uniaxial compressive strength or crack initiation stress).
    \item If $\sigma_{\theta\theta}$ exceeds the strength, purely elastic analysis is insufficient and plasticity must be considered.
    \item Evaluate the influence of the stress ratio $K$ on the distribution of tangential stress around the opening.
\end{itemize}

\subsubsection{Illustrative Example}
Consider a tunnel at depth $z = 700$ m with unit weight $\gamma = 0.027$ MPa/m, giving a vertical stress $p = 18.9$ MPa. For $K = 0.5$, the crown tangential stress is $\sigma_{\theta\theta} = 18.9 \times (3 \times 0.5 - 1) = 9.45$ MPa (compressive). For $K = 2.0$, $\sigma_{\theta\theta} = 18.9 \times (3 \times 2 - 1) = 94.5$ MPa (compressive).

\textbf{Note:} This calculation assumes a homogeneous, intact elastic medium. It does not account for the presence of joints, fractures, or non-elastic behaviour. A definitive stability assessment requires a failure criterion and an appropriate rock mass strength.

\subsection{Elastic--Plastic Zones -- Hoek (1998)}

\subsubsection{Rationale for Plasticity Analysis}
When the induced stresses exceed the rock mass strength, a plastic zone develops around the excavation. The axisymmetric solution for a circular tunnel in a hydrostatic stress field $p_0$, employing a Mohr--Coulomb failure criterion, is presented here.

\subsubsection{Key Parameters (from Week 4: Mohr--Coulomb)}
\[
\sigma_{cm,MC} = \frac{2c\cos\phi}{1-\sin\phi}, \qquad
k = \frac{1+\sin\phi}{1-\sin\phi},
\]
where $c$ and $\phi$ are the Mohr--Coulomb cohesion and friction angle of the rock mass.

\subsubsection{Critical Internal Pressure}
The internal pressure below which a plastic zone forms is given by:
\[
p_{cr} = \frac{2p_0 - \sigma_{cm,MC}}{1 + k}.
\]
For $p_i > p_{cr}$ the rock mass remains elastic; for $p_i < p_{cr}$ a plastic zone of radius $R_p$ develops.

\subsubsection{Plastic Zone Radius -- Unsupported Case ($p_i = 0$)}
The radius of the plastic zone for an unsupported tunnel is:
\[
R_p = a \left[ \frac{2\left(p_0(k - 1) + \sigma_{cm,MC}\right)}{(1 + k)\sigma_{cm,MC}} \right]^{\frac{1}{k-1}}.
\]

\subsubsection{Stresses in the Plastic Zone -- Unsupported Case}
Within the plastic zone ($a \leq r \leq R_p$) and for $p_i = 0$, the radial and tangential stresses are:
\[
\sigma_r^{\text{plastic}} = \frac{\sigma_{cm,MC}}{k - 1}\left[\left(\frac{r}{a}\right)^{k-1} - 1\right],
\]
\[
\sigma_\theta^{\text{plastic}} = k\,\sigma_r^{\text{plastic}} + \sigma_{cm,MC}.
\]

\textbf{Important:} These expressions are valid only for the unsupported condition. For supported tunnels ($p_i > 0$) different formulations apply; these can be found in specialised literature or implemented numerically.

\subsection{Convergence--Confinement Method (CCM)}

\subsubsection{Fundamental Principle}
The CCM conceptualises tunnel excavation as a progressive reduction of the internal support pressure $p_i$ from the initial in-situ stress $p_0$ down to a lower value, while tracking the corresponding inward displacement $u_i$ of the tunnel wall. The method consists of three interrelated components:
\begin{enumerate}[leftmargin=*]
    \item \textbf{Ground Reaction Curve (GRC):} relationship $p_i$ vs $u_i$ for the rock mass.
    \item \textbf{Support Characteristic Curve (SCC):} relationship $p_i$ vs $u_i$ for the installed support system.
    \item \textbf{Longitudinal Deformation Profile (LDP):} describes how $u_i$ varies with distance from the tunnel face, allowing the displacement at support installation to be determined.
\end{enumerate}

\subsubsection{Ground Reaction Curve (GRC)}
Several analytical formulations for the GRC are available; the choice depends on the expected ground behaviour. Table 1 summarises two formulations presented in the lecture slides.

\begin{table}[h!]
\centering
\begin{tabularx}{\textwidth}{l X}
\toprule
\textbf{Reference} & \textbf{Description} \\
\midrule
Carranza-Torres (2004) -- $GRC_{2004}^{GHB}$ & Small-strain generalised Hoek--Brown model with dilation and residual strength. \\
Vrakas \& Anagnostou (2014) -- $GRC_{2014}^{MC}$ & Large-strain Mohr--Coulomb model with dilation. \\
\bottomrule
\end{tabularx}
\caption{GRC formulations discussed in Week 7.}
\end{table}

\textbf{Model uncertainty:} Different constitutive assumptions can lead to significantly different predictions. Both models should be considered as bounds in sensitivity analyses, consistent with the uncertainty theme introduced in Week 5.

\textbf{Mohr--Coulomb GRC (simplified, no dilation):}

The GRC has two branches, divided by the critical pressure $p_{cr}$:

\textit{Elastic branch} ($p_i \geq p_{cr}$):
\[
u_r = \frac{(1+\nu)(p_0 - p_i)\,a}{E_m}
\]

\begin{notebox}[Why $E_m$ and not $E_i$?]
The modulus appearing in any tunnel deformation formula must represent the material that is actually deforming around the opening. The \textit{rock mass} surrounding a tunnel contains joints, fractures, and discontinuities --- these dominate the bulk deformation response at tunnel scale. The intact rock modulus $E_i$ (measured on a $\sim$50~mm laboratory core) describes only the stiffness between joints and would dramatically underpredict convergence. The reduction is significant: for Snowy 2.0, $E_i = 26$~GPa but $E_m \approx 13.5$~GPa at GSI~=~60 and only $\sim$2.1~GPa at GSI~=~30. By analogy, the SCC uses $E_c$ (concrete modulus) because the lining is what deforms there. Each modulus must match the material whose deformation it controls.
\end{notebox}

\textit{Plastic branch} ($p_i < p_{cr}$): First compute the plastic zone radius at the current $p_i$:
\[
R_p(p_i) = a \left[ \frac{2\left(p_0(k-1)+\sigma_{cm,MC}\right)}{(1+k)\left(p_i(k-1)+\sigma_{cm,MC}\right)} \right]^{\frac{1}{k-1}}
\]
Then the wall displacement (approximate):
\[
u_r \approx u_r(p_{cr}) \times \left(\frac{R_p(p_i)}{a}\right)^2
\]
where $u_r(p_{cr}) = (1+\nu)(p_0-p_{cr})\,a/E_m$ is the elastic displacement at the onset of yielding.

\begin{notebox}[GRC shape]
The GRC starts at $(u_r=0,\; p_i=p_0)$ and ends at $(u_r=u_{\max},\; p_i=0)$. The elastic branch is linear with slope $E_m/[(1+\nu)\,a]$. At $p_i = p_{cr}$ the curve transitions to the plastic branch, which is concave (displacement increases rapidly as $p_i$ decreases). The unsupported case ($p_i = 0$) gives the maximum convergence $u_{\max}$.
\end{notebox}

\subsubsection{Support Characteristic Curve (SCC)}
Key parameters defining the SCC (from the slides) are:
\begin{itemize}
    \item $k_s$: support stiffness (slope of the SCC).
    \item $u_{s0}$: radial displacement at the moment of support installation (obtained from the LDP).
    \item $u_{sm}$: maximum elastic displacement capacity of the support.
    \item $p_{sm}$: maximum support pressure capacity.
\end{itemize}

\textbf{SCC for a thin concrete ring} (segmental lining, thickness $t$, radius $a$):
\[
p_{sm} = \frac{\sigma_{ci,\text{lining}} \cdot t}{a}, \qquad
k_s = \frac{E_c \cdot t}{(1-\nu_c^2)\,a^2}, \qquad
u_{sm} = \frac{p_{sm}}{k_s}
\]

The SCC is a straight line starting at $(u_{s0},\; 0)$ with slope $k_s$ until it reaches $p_{sm}$:
\[
p_i = k_s \cdot (u_r - u_{s0}) \quad \text{for } u_{s0} \leq u_r \leq u_{s0}+u_{sm}; \qquad p_i = p_{sm} \text{ thereafter.}
\]

\begin{warnbox}[SCC position depends on installation timing]
The SCC starts at $u_{s0}$, which is determined by the LDP. Moving $u_{s0}$ to the left (earlier installation) shifts the entire SCC leftward, causing the GRC--SCC intersection to occur at a higher $p_{eq}$ (more load on the support). Moving $u_{s0}$ to the right (later installation) reduces $p_{eq}$ but allows more convergence before support activation.
\end{warnbox}

\subsubsection{Longitudinal Deformation Profile (LDP)}
The LDP provides the radial displacement that has already taken place before support is installed. Two commonly used LDP formulations are:

\textit{Elastic LDP (Panet 1995):}
\[
u^*(X^*) =
\begin{cases}
(1 - \alpha)\,e^{1.5X^*}, & X^* \leq 0 \quad \text{(ahead of face)},\\[4pt]
1 - \alpha\,e^{-1.5X^*}, & X^* > 0 \quad \text{(behind face)},
\end{cases}
\]
with $u^* = u/u_{\max}$ and $X^* = X/a$, where $X$ is distance from the tunnel face (positive behind, negative ahead). The parameter $\alpha$ represents the fraction of total convergence occurring behind the face; Panet's original calibration uses $\alpha = 0.75$, which gives $u^* = 0.25$ at the face ($X^* = 0$) on both branches. Substituting $\alpha = 0.75$ into the behind-face branch recovers the familiar form $u = u_{\max}[0.25 + 0.75(1 - e^{-1.5X^*})]$.

\textit{Elastic-plastic LDP (Vlachopoulos \& Diederichs 2009):}
\[
u^*(X^*) =
\begin{cases}
\frac{1}{3} e^{2X^* - 0.15R^*}, & X^* \leq 0 \quad \text{(ahead of face)},\\[4pt]
1 - \left[1 - \frac{1}{3} e^{-0.15R^*}\right] e^{\frac{-3X^*}{R^*}}, & X^* > 0 \quad \text{(behind face)},
\end{cases}
\]
with $u^* = u/u_{\max}$, $X^* = X/a$, and $R^* = R_p/a$.

\begin{notebox}[How to compute $u_{s0}$ from the LDP]
Given installation distance $L$ behind the face:
\begin{enumerate}[nosep]
\item Compute $X^* = L/a$ and $R^* = R_p/a$ from the GRC.
\item Evaluate $u^*$ using the V\&D formula (the $X^*>0$ branch for installation behind the face).
\item $u_{s0} = u^* \times u_{\max}$, where $u_{\max}$ is the GRC displacement at $p_i = 0$.
\end{enumerate}
\end{notebox}

\subsubsection{Influence of Installation Timing on $u_{s0}$}
The LDP yields $u_{s0}$, the radial displacement at the location where support is installed. This value determines the starting point of the SCC on the displacement axis. The relationship between installation timing and overall stability is non-monotonic and requires careful optimisation:

\begin{itemize}
    \item \textbf{Earlier installation (smaller $u_{s0}$):} The support engages while the rock mass retains a greater portion of its load-bearing capacity, thereby limiting ultimate convergence. However, if the ground continues to deform significantly, the support may be subjected to higher loads. The net benefit depends on the stiffness and capacity of the support relative to the GRC.
    \item \textbf{Later installation (larger $u_{s0}$):} A larger proportion of the total convergence occurs before support activation, reducing the load transferred to the support. Nevertheless, excessive pre-convergence may compromise tunnel functionality or lead to instability prior to support installation.
\end{itemize}

The optimal installation point is determined through iterative CCM analysis, balancing these competing effects to achieve an acceptable factor of safety while maintaining serviceability limits.

\subsubsection{Equilibrium and Factor of Safety}
The intersection of the GRC and SCC defines the equilibrium state, characterised by the equilibrium pressure $p_{eq}$ and displacement $u_{eq}$. The factor of safety is then computed as:
\[
\text{FoS} = \frac{p_{sm}}{p_{eq}}.
\]

\textbf{Note on FoS targets:} Acceptable values of the factor of safety are project-specific and depend on asset criticality, design life, consequences of failure, uncertainty in input parameters, and the use of observational methods. Typical targets for permanent tunnels range from 1.5 to 2.0, but engineering judgment and project specifications must always be consulted.

\begin{notebox}[Is RocSupport essential to compute the FoS?]
\textbf{No.} The entire CCM calculation --- including finding $p_{eq}$ --- can be performed in a spreadsheet (Excel or equivalent). RocSupport is convenient because it automates the GRC--SCC intersection and provides built-in plotting, but it is not doing anything the user cannot do manually.

\vspace{0.5em}


\textbf{Finding $p_{eq}$ in Excel:} The intersection is a single nonlinear equation. Two practical approaches:
\begin{itemize}[nosep]
    \item \textbf{Goal Seek (recommended):} Tabulate $p_i$ from $p_0$ down to $0$, compute the corresponding GRC displacement $u_r$ and the SCC pressure at the same displacement $k_s(u_r - u_{s0})$, and use Goal Seek to find the $p_i$ where $p_{\mathrm{GRC}}(u_r) = p_{\mathrm{SCC}}(u_r)$.
    \item \textbf{Graphical inspection:} Plot GRC and SCC on the same $(u_r, p_i)$ axes and read off the intersection. Precision is lower but adequate for a design check.
\end{itemize}
\vspace{0.5em}

\textbf{Where RocSupport (or equivalent) helps:} non-circular tunnels, multiple support layers, complex support geometries (e.g., steel sets + shotcrete + bolts), residual-strength formulations, and parametric sensitivity studies. For the axisymmetric, single-lining problems addressed in this unit, a well-constructed spreadsheet gives the same answer.
\end{notebox}

\subsection{Case Study: Gotthard Motorway Tunnel}
\textbf{Input parameters:} $a = 6$ m, depth $60$ m, $p_0 = 1.62$ MPa, intact UCS $= 7$ MPa, GSI $= 15$, $m_i = 10$, $D = 0$.

\textbf{Rock mass properties (Week 4):} $\sigma_{cm} = 0.47$ MPa, $E_m = 353$ MPa, $\nu = 0.3$, $\phi = 25.64^\circ$.

\subsubsection{Comparison of GRC Solutions}
Table 2 compares predictions obtained with two different GRC formulations.

\begin{table}[h!]
\centering
\begin{tabular}{lcc}
\toprule
\textbf{Parameter} & \textbf{$GRC_{2004}^{GHB}$} & \textbf{$GRC_{2014}^{MC}$} \\
\midrule
Plastic zone radius (unsupported) & 22.4 m & 13.5 m \\
Tunnel convergence (unsupported) & 4.99\% & 1.97\% \\
\bottomrule
\end{tabular}
\caption{Comparison of plastic zone radii and tunnel convergence predicted by two GRC formulations (from Week 7 slides).}
\end{table}

The substantial difference highlights the influence of model uncertainty; both solutions should be considered as plausible bounds in design.

\subsubsection{Effect of Installation Timing ($GRC_{2014}^{MC}$, $LDP_{2009}^{EP}$)}
Table 3 illustrates how the factor of safety varies with the timing of support installation.

\begin{table}[h!]
\centering
\begin{tabular}{lccc}
\toprule
\textbf{Installation criterion} & $L = 3$ m & $u_{s0} = 40$ mm & $\epsilon = 0.3\%$ \\
\midrule
Resulting $u_{s0}$ (mm) & 53.7 & 40.0 & 18.0 \\
$R_p$ supported (m) & 10.0 & 9.2 & 8.28 \\
Convergence (\%) & 1.0 & 0.8 & 0.64 \\
FoS & 1.93 & 1.38 & 0.82 \\
\bottomrule
\end{tabular}
\caption{Influence of support installation timing on plastic zone radius, tunnel convergence, and factor of safety (from Week 7 slides). The three criteria correspond to installing the support at $L=3$~m behind the face, when convergence reaches $u_{s0}=40$~mm, and when the strain reaches $\epsilon=0.3\%$, respectively.}
\end{table}

As installation occurs earlier (i.e., $u_{s0}$ decreases from 53.7~mm to 40~mm to 18~mm), the factor of safety correspondingly declines from 1.93 to 0.82. This trend demonstrates that earlier support engagement increases the load demand on the support system, potentially leading to an unsafe design if installation is too early.

\subsection{Snowy 2.0 Worked Example (650 m depth)}

Using the Snowy 2.0 data ($\sigma_{ci}=72$~MPa, $E_i=26$~GPa, $\nu = 0.22$, $\gamma=26$~kN/m$^3$, $a=4.75$~m, $L=3$~m):

\begin{notebox}[Rock mass and lining assumptions for this worked example]
\textbf{Rock mass:} The Snowy 2.0 rock mass consists of interbedded sandstone and siltstone. To reflect the influence of weaker layers and provide a conservative estimate, $m_i = 7$ (siltstone) is adopted.
\vspace{0.5em}

\textbf{Concrete lining (illustrative only):} $\sigma_{ci,\text{lining}} = 60$~MPa, $E_c = 40$~GPa, $\nu_c = 0.25$, thickness $t = 0.5$~m.
\vspace{0.5em}

\textbf{Note for Assignment 2:} Students should use $m_i = 17$ (sandstone) and the lining properties specified in the assignment brief. The values used here are deliberately different to avoid pre-empting the assignment.
\end{notebox}

\subsubsection{Calculation Workflow for Table 4}

The values in Table~4 follow from the workflow below. All inputs above; $p_0 = \gamma z = 26 \times 650/1000 = 16.9$~MPa.

\begin{enumerate}[leftmargin=*,itemsep=2pt]
    \item \textbf{Hoek--Brown parameters} ($D=0$): compute $m_b$, $s$, and $a_{HB}$ from GSI and $m_i$.
    \item \textbf{Rock mass UCS:} $\sigma_{cm} = \sigma_{ci} \cdot s^{a_{HB}}$.
    \item \textbf{Rock mass modulus} (Hoek--Diederichs 2006): $E_m = E_i \,[0.02 + 1/(1 + e^{(60-GSI)/11})]$.
    \item \textbf{Equivalent Mohr--Coulomb} ($\phi$, $c$): obtained by fitting the H--B envelope using RocData. For deep tunnels where the plastic zone extends well into the rock mass, the fit stress range is typically taken up to $\sigma_{3,\max} \approx p_0$ so that the linearised MC envelope spans stresses from the tunnel wall ($\sigma_3 \approx 0$) out to the elastic boundary ($\sigma_3 \approx p_0$). Hoek (2002) provides an iterative formula relating $\sigma_{3,\max}$ to tunnel geometry and in-situ stress; for the Snowy 2.0 example the fit gives the ($\phi$, $c$) values listed in Table~4.
    \item \textbf{MC-derived quantities:} $k$, $\sigma_{cm,MC}$, and $p_{cr}$ from the equations in Sections~1.5.2 and 1.5.3.
    \item \textbf{Plastic zone radius:} apply the $R_p$ equation in Section~1.5.4 with $p_i = 0$.
    \item \textbf{Maximum convergence:} $u_r(p_{cr}) = (1+\nu)(p_0-p_{cr})a/E_m$, then $u_{\max} \approx u_r(p_{cr}) \cdot (R_p/a)^2$.
    \item \textbf{LDP at $L = 3$~m:} compute $X^* = L/a$, $R^* = R_p/a$, evaluate the V\&D (2009) behind-face branch, and $u_{s0} = u^* \cdot u_{\max}$.
    \item \textbf{Support properties} (independent of GSI): apply the SCC equations in Section~1.6.3 with the lining properties above.
\end{enumerate}

\subsubsection{Results Summary}

\begin{table}[h!]
\centering
\small
\begin{tabular}{lcc}
\toprule
\textbf{Parameter} & \textbf{GSI = 60} & \textbf{GSI = 30} \\
\midrule
$m_b$ & 1.68 & 0.575 \\
$s$ & 0.0117 & 0.000418 \\
$a_{HB}$ & 0.503 & 0.522 \\
\midrule
$\sigma_{cm}$ (MPa) & 7.70 & 1.24 \\
$E_m$ (GPa) & 13.52 & 2.14 \\
$\phi$ (\textdegree) & 31.5 & 23.8 \\
$c$ (MPa) & 2.85 & 1.72 \\
\midrule
$p_0$ (MPa) & 16.90 & 16.90 \\
$\sigma_{cm,MC}$ (MPa) & 10.18 & 5.28 \\
$k$ & 3.19 & 2.35 \\
$p_{cr}$ (MPa) & 5.64 & 8.51 \\
Response & Plastic & Plastic \\
\midrule
$R_p$ (m) & 6.83 & 11.17 \\
$R_p/a$ & 1.44 & 2.35 \\
$u_{\max}$ (mm) & 9.97 & 125.6 \\
$u_{s0}$ at $L=3$~m (mm) & 8.0 & 82.6 \\
\midrule
$p_{sm}$ (MPa) & 6.32 & 6.32 \\
$k_s$ (MPa/m) & 945.6 & 945.6 \\
$u_{sm}$ (mm) & 6.68 & 6.68 \\
\bottomrule
\end{tabular}
\caption{Snowy 2.0 CCM summary using $m_i = 7$ (siltstone-controlled rock mass) and an illustrative 0.5~m thick concrete lining ($\sigma_{ci,\text{lining}}=60$~MPa, $E_c=40$~GPa, $\nu_c=0.25$). Lining properties are independent of GSI.}
\end{table}

\begin{warnbox}[Effect of $m_i$ selection]
Using a lower $m_i$ (siltstone instead of sandstone) significantly reduces $m_b$, leading to lower rock mass strength, larger plastic zones, and increased tunnel convergence. This demonstrates the strong sensitivity of Hoek--Brown parameters to lithological assumptions and highlights the importance of selecting representative values for mixed ground conditions.
\end{warnbox}

\begin{warnbox}[GSI = 30 — solving the GRC--SCC intersection]
A common mistake is to compare the lining's elastic capacity $u_{sm}$ against the \textit{unsupported} remaining convergence $u_{\max} - u_{s0} \approx 43$~mm and conclude the lining fails. This reasoning is incorrect: once the lining engages, it applies pressure on the rock mass, which reduces $R_p$ and therefore reduces the final convergence $u_{eq}$ well below $u_{\max}$. The CCM methodology requires solving the GRC--SCC intersection explicitly.\\[4pt]

\textbf{Intersection analysis for GSI = 30:} For $p_i > 0$, the supported GRC displacement is $u_r(p_i) = u_r(p_{cr}) \cdot (R_p(p_i)/a)^2$ where $R_p(p_i)$ shrinks as $p_i$ increases. Setting this equal to the SCC $p_i = k_s(u_r - u_{s0})$ and iterating (Excel Goal Seek or equivalent) gives $p_{eq} \approx 1.23$~MPa at $u_{eq} \approx 83.9$~mm (just $\sim$1.3~mm of post-installation displacement, well within $u_{sm} = 6.7$~mm). The corresponding factor of safety is:
\[
\mathrm{FoS} = \frac{p_{sm}}{p_{eq}} = \frac{6.32}{1.23} \approx 5.1
\]
A similar analysis for GSI = 60 gives $p_{eq} \approx 0.71$~MPa, FoS $\approx 8.9$. So the lining is nominally adequate for both conditions in this idealised CCM model.

\textbf{Caveats — why numerical modelling (Week 8) is still needed:} The CCM result above is optimistic because the analytical framework has limitations that are acute in poor ground:
\begin{itemize}[nosep]
    \item \textbf{Axisymmetric assumption} ignores gravity asymmetry, non-circular deformations, and stress-induced anisotropy that become significant at large $R_p/a$.
    \item \textbf{Pre-convergence control:} $u_{s0} \approx 83$~mm ahead of the lining is a very large strain ($\sim$1.7\% of diameter). In reality, the tunnel face may lose stability before this displacement can develop, requiring face bolts, forepoling, or probe drilling.
    \item \textbf{Time-dependent effects} (creep, squeezing) are not captured and can double long-term convergence.
    \item \textbf{Local failures} (block detachment, ravelling, structurally-controlled wedges) at $R_p/a > 2$ are outside CCM's scope.
\end{itemize}

\textbf{Practical design implications:} Even though the simplified CCM gives FoS $\sim$5 for GSI = 30, experienced designers would add contingency measures for poor rock:
\begin{itemize}[nosep]
    \item \textbf{Yielding support elements} (compressible blocks, sliding steel sets) to handle unexpected squeezing;
    \item \textbf{Pre-reinforcement} (rock bolts, forepoling, face dowels) to control the $\sim$83~mm pre-convergence;
    \item \textbf{Observational method} with convergence monitoring to validate the CCM predictions.
\end{itemize}
Use \textbf{RocSupport} (or Excel Goal Seek) to find $p_{eq}$ precisely and compute the FoS. For comprehensive verification including asymmetric loading, non-circular geometry, and construction sequencing, numerical modelling (Week~8) is essential.
\end{warnbox}

\subsection{Assignment 2 Tasks (Week 7 Component)}
The Week~7 component uses the same Snowy 2.0 project data as Week~6 (UCS $= 72$~MPa, $E_i = 26$~GPa, $\nu = 0.22$, $\gamma = 26$~kN/m$^3$). Two rock mass conditions: GSI $= 60$ (good) and GSI $= 30$ (poor).

\begin{enumerate}
    \item \textbf{Kirsch elastic solutions for TBM tunnels:} Calculate $\sigma_\theta$ at crown/invert and sidewalls at various depths (100--800~m). Identify critical regions and how stress concentration changes with depth.

    \item \textbf{Rock mass and elastoplastic parameters for the main access tunnel:} At 650~m depth under both GSI conditions, determine whether the response is elastic or elastoplastic. Calculate $p_{cr}$, elastic wall displacement, and $R_p$. Derive Hoek-Brown parameters ($m_b$, $s$, $a$), $\sigma_{cm}$, $\sigma_m$, $\sigma_{tm}$, and $E_m$.

    \item \textbf{CCM and RocSupport analysis:} Assess ground-support interaction for the main access tunnel using the CCM and \textbf{RocSupport software}. Use the lining properties specified in the assignment brief, installed at $L = 3$~m from the face. Determine plastic radius, equilibrium pressure, FoS, and support performance. Include GRC, SCC, and LDP curves.

    \item \textbf{Brief reflection:} Comment on the usefulness of Kirsch solutions, differences between elastic and elastoplastic interpretations, CCM/RocSupport for support assessment, and issues requiring numerical verification.
\end{enumerate}

\begin{notebox}[Software submission]
Submit RocSupport \texttt{.rsp} files with your report. Keep all formulas visible in calculation spreadsheets.
\end{notebox}

\subsection{Week 7 Summary Table}
\begin{table}[h!]
\centering
\begin{tabularx}{\textwidth}{l X X}
\toprule
\textbf{Component} & \textbf{Purpose} & \textbf{Key assumptions / cautions} \\
\midrule
Kirsch (1898) & Elastic boundary stress for circular tunnel & Elastic, homogeneous, isotropic; invalid once yielding occurs. \\
Hoek (1998) & Plastic radius $R_p$ and critical pressure $p_{cr}$ & Circular geometry, hydrostatic stress, Mohr--Coulomb idealisation. \\
GRC & Relationship between support pressure and convergence & Choice of constitutive model strongly affects results. \\
SCC & Relationship between support pressure and displacement & Requires realistic stiffness, capacity, and $u_{s0}$. \\
LDP & Displacement profile ahead of/behind the face & Model selection (elastic vs.\ elastic-plastic) influences $u_{s0}$. \\
Equilibrium $p_{eq}$ & Intersection of GRC and SCC & Sensitive to all input parameters and assumptions. \\
FoS ($p_{sm}/p_{eq}$) & Capacity-to-demand ratio & Target values are project-specific; engineering judgment essential. \\
\bottomrule
\end{tabularx}
\caption{Summary of Week 7 concepts, their purpose, and associated caveats.}
\end{table}

\vspace{1cm}
\noindent\textit{End of Week 7 Study Notes}

\end{document}